\usepackage[a4paper,margin=2.6cm,headheight=14pt]{geometry}
\usepackage{amsmath,amssymb,mathtools}
\usepackage{amsthm}
\usepackage{enumitem}
\usepackage{xcolor}
\usepackage{microtype}
\usepackage[unicode,colorlinks=true,linkcolor=blue!55!black,urlcolor=blue!55!black]{hyperref}
\usepackage[nameinlink,noabbrev]{cleveref}

\ctexset{
  chapter = {
    name = {第,讲},
    number = \arabic{chapter},
    aftername = \quad
  }
}

\setcounter{tocdepth}{2}
\setcounter{secnumdepth}{2}

% 定理类环境共享同一计数器，并按讲次重新编号。
\theoremstyle{plain}
\newtheorem{theorem}{定理}[chapter]
\newtheorem{lemma}[theorem]{引理}
\newtheorem{corollary}[theorem]{推论}
\newtheorem{proposition}[theorem]{命题}

\theoremstyle{definition}
\newtheorem{definition}[theorem]{定义}
\newtheorem{example}[theorem]{例}

\theoremstyle{remark}
\newtheorem*{remark}{注}
\newtheorem*{question}{思考}

\crefname{theorem}{定理}{定理}
\crefname{lemma}{引理}{引理}
\crefname{corollary}{推论}{推论}
\crefname{proposition}{命题}{命题}
\crefname{definition}{定义}{定义}
\crefname{example}{例}{例}
\crefname{equation}{式}{式}
\crefname{section}{节}{节}
\crefformat{chapter}{第#2#1#3讲}

\newcommand{\bits}{\{0,1\}}
\newcommand{\bitsplus}{\bits^{+}}
\newcommand{\bitsstar}{\bits^{*}}
\newcommand{\eps}{\varepsilon}
\newcommand{\card}[1]{\lvert #1\rvert}
\newcommand{\Nat}{\mathbb{N}}

\renewcommand{\proofname}{证明}
